% =============================================================================
% BLOC 4 — 10 slides / 20 minutes
% =============================================================================
\section{Bloc 4 — Piloter}

% 34
\sectiondivider{BLOC 4}{Piloter l'équipe du projet}{de développement logiciel}{« Cinq personnes, neuf mois, un client réel — et une capacité réelle qui chute au moment où il faut accélérer. » · C4.1 à C4.6}

% 35 — anglais
\begin{frame}{Leading a team whose real capacity changed}
\eyebrow{In English · 3 to 4 minutes}
\selectlanguage{english}
\begin{center}
\begin{tikzpicture}[node distance=.35cm,every node/.style={font=\small,align=center}]
\node[draw=Line,fill=white,rounded corners,minimum width=2.65cm,minimum height=1.35cm] (a) {\textbf{The plan}\\M2 frame the project\\M1 scale development};
\node[draw=Red,fill=SoftRed,rounded corners,minimum width=2.65cm,minimum height=1.35cm,right=of a] (b) {\textbf{The shock}\\Another M1 project\\is postponed};
\node[draw=Orange,fill=SoftOrange,rounded corners,minimum width=2.65cm,minimum height=1.35cm,right=of b] (c) {\textbf{The consequence}\\Late arrival $\neq$\\instant capacity};
\node[draw=Blue,fill=SoftBlue,rounded corners,minimum width=2.65cm,minimum height=1.35cm,right=of c] (d) {\textbf{My response}\\Reassign · simplify\\postpone · protect};
\draw[-{Latex},Blue,thick] (a)--(b);\draw[-{Latex},Blue,thick] (b)--(c);\draw[-{Latex},Blue,thick] (c)--(d);
\end{tikzpicture}
\end{center}
\vspace{0.7em}
\begin{columns}[T]
\begin{column}{0.48\linewidth}
\begin{block}{What I learned}
\small Capacity must be reviewed like a risk: regularly, by person and by domain.
\end{block}
\end{column}
\begin{column}{0.48\linewidth}
\begin{block}{Limits of the experience}
\small No recruitment, no HR process, no formal career management. I separate what I did from what I would implement professionally.
\end{block}
\end{column}
\end{columns}
\selectlanguage{french}
\end{frame}

% 36
\begin{frame}{Le cadre, et ce qu'il autorise ou non}
\eyebrow{Bloc 4 · Périmètre de la preuve}
\begin{columns}[T]
\begin{column}{0.32\linewidth}
\begin{block}{FAIT}
\small
\begin{tightitems}
  \item Coordination de 5 personnes
  \item Rôles et responsabilités
  \item Arbitrages de périmètre
  \item Rétrospective + entretiens
\end{tightitems}
\end{block}
\end{column}
\begin{column}{0.32\linewidth}
\begin{block}{NON RÉALISÉ}
\small
\begin{tightitems}
  \item Recrutement
  \item Processus RH
  \item Référentiel individuel formel
  \item Plans de carrière
\end{tightitems}
\end{block}
\end{column}
\begin{column}{0.32\linewidth}
\begin{block}{EN ENTREPRISE}
\small
\begin{tightitems}
  \item Fiches de rôle
  \item Mobilité / formation / recrutement
  \item Référentiel collectif + individuel
  \item Feedback continu
\end{tightitems}
\end{block}
\end{column}
\end{columns}
\vspace{0.5em}
\begin{block}{Positionnement}
\centering\small \strong{Ce que je peux démontrer, c'est ce que j'ai réellement piloté ; ce que je ne prétends pas, je le transforme en dispositif cible.}
\end{block}
\end{frame}

% 37
\begin{frame}{Raisonner par domaine, pas par personne}
\eyebrow{C4.1 · Compétences nécessaires}\competency{C4.1}
\small
\begin{center}
\begin{tabularx}{0.96\linewidth}{*{4}{>{\centering\arraybackslash}X}}
\strong{Pilotage produit}\\\scriptsize périmètre · client · risques & \strong{Mobile Flutter}\\\scriptsize parcours apprenant & \strong{Web admin}\\\scriptsize gestion métier & \strong{Backend + données}\\\scriptsize auth · RLS · migrations \\[1.0em]
\strong{UX / UI}\\\scriptsize parcours · accessibilité & \strong{Tests + sécurité}\\\scriptsize non-régression & \strong{Paiements}\\\scriptsize cycle d'abonnement & \strong{Déploiement + doc}\\\scriptsize reprise · transmission
\end{tabularx}
\end{center}
\vspace{0.75em}
\begin{block}{Règle}
\centering\small Chaque domaine doit avoir \strong{un responsable identifiable et une solution de continuité}. Une compétence disponible sur le papier n'est pas une capacité utilisable si la personne n'est pas disponible.
\end{block}
\end{frame}

% 38
\begin{frame}{Constituer une équipe sans la recruter}
\eyebrow{C4.2 · Constitution de l'équipe}\competency{C4.2}
\begin{columns}[T]
\begin{column}{0.49\linewidth}
\begin{block}{Ce qui a été fait}
\small
\begin{tightitems}
  \item Groupe imposé par le cadre pédagogique
  \item Rôles précisés : chef de projet, PO, Scrum Master, développement
  \item RACI sur 10 activités
  \item Circuits de validation et continuité
\end{tightitems}
\end{block}
\vspace{0.25em}
\small\warn{La RACI a révélé un cumul de rôles sur deux personnes : efficace pour la continuité, risqué pour la dépendance.}
\end{column}
\begin{column}{0.47\linewidth}
\begin{block}{Plan de constitution cible}
\small Une fiche par rôle :
\begin{tightitems}
  \item Mission + responsabilités
  \item Compétences techniques
  \item Compétences relationnelles
  \item Autonomie recherchée
  \item Interactions clés
\end{tightitems}
\end{block}
\vspace{0.25em}
\small\strong{Arbitrage ensuite : mobilité interne, formation ciblée ou recrutement externe.}
\end{column}
\end{columns}
\end{frame}

% 39
\begin{frame}{Coordonner quand la capacité s'effondre}
\eyebrow{C4.3 · Gestion opérationnelle}\competency{C4.3}
\begin{center}
\begin{tikzpicture}[node distance=.25cm,every node/.style={font=\scriptsize,align=center}]
\node[draw=Line,fill=white,rounded corners,minimum width=2.65cm,minimum height=1.55cm] (a) {\textbf{1 · Le plan}\\M2 cadrage + transverse\\M1 développement};
\node[draw=Red,fill=SoftRed,rounded corners,minimum width=2.65cm,minimum height=1.55cm,right=of a] (b) {\textbf{2 · Le choc}\\Projet DEV M1 décalé\\au pire moment};
\node[draw=Orange,fill=SoftOrange,rounded corners,minimum width=2.65cm,minimum height=1.55cm,right=of b] (c) {\textbf{3 · Conséquence}\\Onboarding tardif\\= capacité différée};
\node[draw=Blue,fill=SoftBlue,rounded corners,minimum width=2.65cm,minimum height=1.55cm,right=of c] (d) {\textbf{4 · Réponse}\\Réaffecter · simplifier\\reporter · protéger};
\draw[-{Latex},Blue,thick] (a)--(b);\draw[-{Latex},Blue,thick] (b)--(c);\draw[-{Latex},Blue,thick] (c)--(d);
\end{tikzpicture}
\end{center}
\vspace{0.75em}
\begin{block}{Leçon}
\centering\small Le pilotage consiste aussi à décider \strong{ce qu'on cesse de suivre ou de construire} quand la capacité réelle change.
\end{block}
\end{frame}

% 40
\begin{frame}{Accompagner, et le coût réel de la transmission}
\eyebrow{C4.4 + C4.5 · Animation et montée en compétences}\competency{C4.4 · C4.5}
\begin{columns}[T]
\begin{column}{0.48\linewidth}
\begin{block}{Accompagnement quotidien}
\small Points d'équipe · Teams · Confluence · Jira · GitHub · aide technique ciblée.
\end{block}
\begin{block}{Montée en compétences}
\small Documentation officielle → prototype limité → tâche réelle.
\end{block}
\end{column}
\begin{column}{0.48\linewidth}
\begin{block}{La limite structurelle}
\small Former quelqu'un coûte plus de temps à court terme que confier la tâche à la personne déjà experte.
\end{block}
\begin{block}{Ce que je ferais}
\small Binômage · revues techniques · démonstrations internes · rotation sur certains domaines.
\end{block}
\end{column}
\end{columns}
\vspace{0.55em}
\begin{block}{Point clé}
\centering\small La documentation réduit la dépendance, mais \strong{elle ne remplace pas le transfert actif de connaissance}.
\end{block}
\end{frame}

% 41
\begin{frame}{Évaluer : ce que j'ai fait, ce que je construirais}
\eyebrow{C4.6 · Suivi de la performance}\competency{C4.6}
\begin{columns}[T]
\begin{column}{0.44\linewidth}
\begin{block}{Réalisé}
\small
\begin{tightitems}
  \item Jalons, exigences, risques, capacité
  \item Rétrospective générale
  \item Entretiens individuels
  \item Retour mutuel, pas seulement descendant
\end{tightitems}
\end{block}
\small\strong{La rétrospective a montré ce que Jira ne voit pas : utilité ressentie, autonomie, clarté du rôle, charge perçue.}
\end{column}
\begin{column}{0.52\linewidth}
\scriptsize
\begin{tabularx}{\linewidth}{>{\bfseries}p{1.8cm}XX}
\toprule
 & Collectif & Individuel \\
\midrule
Livraison & objectifs atteints & engagements tenus \\
Qualité & non-régression · revues & qualité des livrables \\
Collaboration & blocages partagés & communication · feedback \\
Autonomie & dépendances réduites & avancer + expliciter \\
Progression & besoins identifiés & compétences du rôle \\
\bottomrule
\end{tabularx}
\vspace{0.35em}
\small\strong{Cible : feedbacks courts tous les 2–3 sprints, à double sens.}
\end{column}
\end{columns}
\end{frame}

% 42
\begin{frame}{Préparer la reprise avant la fin du projet}
\eyebrow{C3.6 + C4.5 · Transmission}\competency{C3.6}
\begin{columns}[T]
\begin{column}{0.52\linewidth}
\begin{block}{Ce qui a été produit}
\small
\begin{tightitems}
  \item MkDocs développeur > 100 pages
  \item Monorepo structuré
  \item DevContainer
  \item Migrations + scripts d'initialisation
  \item Environnements séparés
\end{tightitems}
\end{block}
\centering\repofigure[0.23\textheight]{figures/ipf/acceuil_doc_dev_mkdocs.png}
\end{column}
\begin{column}{0.44\linewidth}
\begin{block}{La transmissibilité ne tient pas à un document}
\small Le projet doit pouvoir être installé et exécuté ailleurs que sur les postes de ceux qui l'ont développé.
\end{block}
\begin{block}{Transmission client}
\small Une réunion de reprise est prévue pour présenter l'état réel du produit, les limites, les procédures de lancement et les conditions de reprise par une future équipe.
\end{block}
\vspace{0.25em}
\scriptsize\muted{Le compte rendu ne sera ajouté à l'oral que si cette réunion a effectivement eu lieu avant la soutenance.}
\end{column}
\end{columns}
\end{frame}

% 43
\begin{frame}{Bilan critique du bloc 4}
\eyebrow{Prise de recul}
\small
\begin{tabularx}{\linewidth}{>{\bfseries\color{Blue}}p{3.0cm}X}
\toprule
Sur la planification & Un planning qui glisse reste utile s'il rend l'écart visible et force l'arbitrage. \\
\addlinespace
Sur les risques & Un risque n'a de valeur que s'il est réévalué et déclenche une décision assez tôt. \\
\addlinespace
Sur la délégation & Le choix le plus rapide à court terme peut concentrer la connaissance et coûter cher à la reprise. \\
\addlinespace
Sur le cumul des rôles & Chef de projet + référent technique donne une vision globale, mais crée un point unique de décision et d'escalade. \\
\bottomrule
\end{tabularx}
\vspace{0.65em}
\begin{block}{À retenir}
\centering\small Le produit final est une base fonctionnelle, versionnée et documentée ; mon principal apprentissage porte sur \strong{la capacité réelle, la délégation et la continuité du collectif}.
\end{block}
\end{frame}
